\section{Introduction}\label{sec:introduction}

Standard survival models assign vanishing probability to remaining event
free as follow-up grows arbitrarily long. That limit is inappropriate
when a clinically meaningful subgroup will never experience the event of
interest, a situation indicated by a sustained plateau in the survival
curve under mature follow-up. Cure models describe this setting by
separating the long-term event-free probability from the timing of
events among individuals who remain susceptible, which makes the cured
proportion an interpretable quantity in oncology and other long-term-survivor
studies \cite{berkson1952,tsodikov2003}.

{
Accommodating this plateau requires a model that allows a positive
limiting survival probability, and three main constructions are available
in the literature for this purpose, differing in the mechanism by which
they generate the cure mass. The first, the mixture cure model,
introduces a latent binary indicator that separates long-term event-free
individuals from those who remain susceptible, so that the cured
proportion and the survival of the susceptible subpopulation are
specified as separate components, typically a regression model for
incidence together with a proper distribution for the latency
\cite{boag1949,berkson1952,farewell1982,sy2000}. The second, the
defective-distribution model, does not represent the cured fraction
through a separate component; instead, the cured fraction arises
intrinsically from the parameter space of a single parametric survival
distribution. For most parameter values the distribution is proper and
the survival function decays to zero, as in conventional survival models;
over a distinct region of the same parameter space the distribution
becomes improper and the survival function instead converges to a
strictly positive limit, which is the cured fraction. A single family
then governs both the failure times of the susceptible individuals and
the size of the cured group, with no separate incidence model and with
covariates entering through its parameters
\cite{balka2009,rocha2017,scudilio2019}. The third, based on latent
competing risks, generates the cure mass from an unobserved count of
competing causes, so that the cured fraction is The probability that this
count is zero \cite{yakovlev1996,chen1999,tsodikov2003,rodrigues2009unif}.
} %This construction can be re-expressed as a mixture whose cured proportion equals that zero-count probability and whose susceptible latency is induced by the count law, so it is a structured special case of the mixture model in which the incidence and latency components are linked through shared parameters rather than specified separately \cite{rodrigues2009}.
{
This mechanism admits a two-stage description. An initiation stage sets
the unobserved number $N$ of latent causes carried by an individual, and
a maturation stage assigns each cause an activation time; the earliest
activation produces the event, so $N=0$ describes an individual who never
fails and identifies the cure mass \cite{rodrigues2025}. Encoding the
count law through its probability-generating function links the latency
distribution to both the population survival and the cure fraction, and
it is this representation on which the model developed here is built.
}

{
In the promotion-time construction the mean number of latent causes is
almost always taken to be log-linear in the covariates
\cite{chen1999,rodrigues2009}. This specification is convenient and
keeps the cure fraction in closed form, but it forces every covariate
to act monotonically on the latent-cause intensity, so threshold effects,
interactions, and periodic or otherwise nonlinear associations cannot be
represented. Because the intensity governs the cure fraction, a
misspecified incidence model biases both the estimated covariate effects
and the cured proportion \cite{amico2018}. The activation-time component
carries another issue: fixing the latency to a restrictive family can be
equally misspecified, and in a cure model this error propagates to the
covariate effects and to the cure fraction. A flexible parametric family
mitigates this while preserving interpretability and a tractable
likelihood, and the generalized gamma is a natural choice, since it
nests the exponential, Weibull, gamma, and log-normal distributions and
spans a wide range of hazard shapes
\cite{stacy1962,prentice1974,cox2007}. Together, these concerns motivate
treating both components flexibly, through a flexible latent-intensity
predictor for the incidence and a sufficiently general parametric latency
for the activation times.
}

{
Flexible predictors have been proposed in cure modelling: smoothers,
support-vector machines, neural networks, and deep architectures have all
been used to relax the usual linear specification. In the promotion-time
setting, Pal and Aselisewine\cite{pal2023} introduced a support-vector
formulation for the incidence component, and Xie and Yu\cite{xie2021}
used a neural network for an unstructured promotion-time component,
while deep promotion-time models and structured nonlinear cure models
further broaden the available high-dimensional architectures\cite{medina2024,feng2026}.
These methods make the covariate effect on the latent-cause intensity
nonlinear, but they differ substantially in how they treat that
intensity, the latency distribution, and the count structure. To our
knowledge, no existing promotion-time approach combines all three within
a single likelihood-based estimation and prediction framework: a
flexible nonlinear predictor for the latent-cause intensity, which shapes
both the cured fraction and the timing of events; a transparent and
sufficiently general parametric latency; and an explicit count-family
treatment of overdispersion.
}

To address this gap, we develop a generalized gamma kernel-machine
promotion-time cure model (GG-KM) that places the log-intensity of
the latent-cause process in a reproducing-kernel Hilbert space (RKHS)
\cite{kimeldorf1971}, while retaining a flexible parametric
generalized gamma latency distribution \cite{stacy1962} and allowing
either Poisson or negative-binomial latent counts. This formulation
separates two sources of modelling flexibility: the kernel predictor
captures potentially complex nonlinear and interaction effects of
covariates on the latent-cause mechanism, whereas the generalized
gamma distribution provides a flexible parametric representation of
the activation-time process.


The remainder of this article is organized as follows. Section~\ref{sec:methods}
presents the proposed generalized gamma kernel-machine promotion-time cure
model, covering the general promotion-time construction, the generalized gamma
latency distribution, the kernel-based latent-cause predictor, and the penalized
EM and ECME estimation procedures with their observed-likelihood refinement and
convergence assessment. Section~\ref{sec:simulation} describes the simulation
study and evaluates the predictive performance of the proposed model against
log-linear, additive-spline, and quadratic alternatives under different
covariate structures and latent-count distributions. Section~\ref{sec:melanoma}
applies the model to the melanoma data and reports its predictive behaviour.
Section~\ref{sec:discussion} discusses the main findings, limitations, and
directions for future research.
